% s-028-p028-m003; original PDF 28, 29
\Aphorism{s-028-p028-m003}{19}
Selbstopferung als Ausdruck eines masslosen Willens zur Macht ist eine häufige Erscheinung in solchen festgezimmerten Gemeinwesen wie etwa die römische Republik war, wo jede Gegenwart ihrer Zukunft und zwar einer langen Zukunft sicher ist, wo eine Erinnerung festgehalten, eine Sitte durchgesetzt, einem Imperativ Gehorsam erzwungen wird über ganze Geschlechterketten hin, wo das Individuum dann, aber auch nur dann für alle Zeiten gelebt hat, wenn es in und mit dem Ganzen gelebt hat. Der Einzelne, der sich inmitten einer solchen Organisation und für sie opfert, übt mit dieser Opferthat einen ungeheuren moralischen Druck aus auf alles, was nach ihm kommt. Er stirbt nicht einsam und unbeachtet, wie die Raupe, die hinter dem Wanderer zertreten im Staube zurückbleibt, seine That geht nicht spurlos verloren, denn sie geschah im Tageslichte des gemeinschaftlichen Bewusstseins und wird, dessen ist er gewiss, dauern solange dieses Bewusstsein dauert. Der Staat wird sie sich bald genug zu eigen machen, als Sporn und Stachel für die Jugend, als Schrecken für den Feind, als Gelöbniss und bindende Formel für die Zukunft, als Schuldabzahlung und verpflichtendes Unterpfand für die staatserhaltenden Götter: so verzichtet der Staatsunterthan auf sein Leben und Glück, um am Leben und Glück des Staates Antheil zu haben, um nach seinem Tode noch als Gespenst, als blosser Name, höchste Machtfülle zu geniessen. Ist dies alles bei uns noch möglich? in einer nicht mehr politomanen Zeit, die nur noch wirthschaftliche Motoren kennt und für die bereits der Romantiker Hegel, mit seiner feierlichen Apotheose der Staatsidee, zu spät kam?
\par

